\chapter{Running a simulation}
\label{ch:runningmd}

Once you have a topology and coordinates, the Molecular Dynamics Manager
(\key{5b}) sets up and runs the simulation. It writes the Amber input files for
each stage, runs them on your computer or writes job scripts for a computing
cluster, shows live progress, and analyzes the results. This chapter covers
setting up and running; Chapter~\ref{ch:analysis} covers analysis.

\begin{glance}
\glanceitem{What it does}{Prepares, runs and monitors a series of simulation
  stages (minimization, heating, equilibration, production).}
\glanceitem{You need it when}{You want to simulate a system built in \key{5a},
  or a \file{.prmtop}/\file{.rst7} pair loaded with \key{1a} \key{2}.}
\glanceitem{Before you start}{A topology and coordinates
  (Chapter~\ref{ch:topology}). An MD engine: \texttt{sander} and
  \texttt{sander.MPI} come with ProPrep; \texttt{pmemd} does not.}
\glanceitem{Where to find it}{\menupath{5b}{Molecular Dynamics Manager}.}
\glanceitem{You get}{A folder \file{simulations/batch_<date>_<time>/} with
  one subfolder per stage, each holding the input, energy output, final
  coordinates and trajectory.}
\end{glance}

\section{The stages of a simulation}\label{sec:runningmd-stages}

A freshly built system cannot simply be switched on: the added hydrogens and
waters clash a little, and every atom is standing still. A \term{protocol}
brings the system gently to the conditions you want, then records it. Most
protocols have four kinds of stage:

\begin{description}
  \item[Minimization] Moves atoms slightly, without time passing, until the
        clashes are gone and the energy stops falling. Measured in
        \emph{cycles}.
  \item[Heating] Gives the atoms speed, raising the temperature from 0 to the
        target (for example 300\,K) over a few picoseconds.
  \item[Equilibration] Lets the system settle. At \term{NPT} (constant
        pressure) the box shrinks or grows until the water has its normal
        density; at \term{NVT} (constant volume) the box size is fixed.
  \item[Production] The run you keep and analyze. Snapshots are written to
        the \term{trajectory} file at regular intervals.
\end{description}

During the early stages, \term{restraints} hold chosen atoms (often the
protein backbone) near their starting positions, so the protein is not
damaged while the water settles around it.

Time in a simulation is counted in \term{steps}. A typical step is 2
femtoseconds (fs). 500 steps make 1\,picosecond (ps), and 1,000\,ps make
1\,nanosecond (ns). Research simulations usually run for tens or hundreds of
nanoseconds.

The settings of each stage live in an Amber input file, the \term{mdin} file.
You rarely need to read it, but a few of its settings are worth knowing:

\begin{tabularx}{\linewidth}{@{}lL@{}}
\toprule
\textbf{Setting} & \textbf{Meaning} \\
\midrule
\texttt{imin=1} & This stage is a minimization (\texttt{imin=0} or absent: dynamics) \\
\texttt{maxcyc}, \texttt{ncyc} & Total minimization cycles; how many of them use the simple ``steepest descent'' method \\
\texttt{nstlim}, \texttt{dt} & Number of steps; length of one step in ps (0.002 = 2\,fs) \\
\texttt{temp0}, \texttt{tempi} & Target and starting temperature (K) \\
\texttt{ntb}, \texttt{ntp} & Constant volume (\texttt{ntb=1}) or constant pressure (\texttt{ntb=2}, \texttt{ntp=1}) \\
\texttt{ntwx}, \texttt{ntpr}, \texttt{ntwr} & Every how many steps to save a snapshot, print energies, write a restart file \\
\texttt{ntr}, \texttt{restraint\_wt}, \texttt{restraintmask} & Restraints on or off, their strength, and which atoms \\
\texttt{\&wt} lines & Changes during the stage, such as a temperature ramp \\
\bottomrule
\end{tabularx}

\section{Step by step: set up and run}\label{sec:runningmd-steps}

The Molecular Dynamics Manager menu:

\begin{screen}
══ Molecular Dynamics Manager Menu ══
  1. ✓ Setup and configure simulations
  2. ○ Execute simulations
       ⚠ Need to setup simulations first
  3. ○ Monitor simulation progress
       ⚠ No simulations currently running
  4. ✓ Analyze completed simulations
  5. ✓ Manage Templates & Protocols
  6. ✓ Manage cluster profiles (HPC configurations)
  7. ✓ Manage run plans (per-protocol resource assignments)
\end{screen}

Setup (\key{1}) is a six-step guide with a progress bar at the top of each
screen:

\begin{screen}
╭─────────────────────────────────────────────────────────────╮
│ Step 1 of 6: Structure Files (Topology + Coordinates)       │
│ [Files] → Protocol → Restraints → Queue → Hardware → Review │
╰─────────────────────────────────────────────────────────────╯
\end{screen}

The route that worked for ubiquitin (shown in full in
Chapter~\ref{ch:tutorial}):

\begin{steps}
  \item \key{5b}, then \key{1}.
  \item \textbf{Files:} check that the right pair is listed, then \key{n}.
  \item \textbf{Protocol:} \key{1} (predefined), \key{4} (Workshop), then
        \key{e} to shorten it or \key{u} to use it as it is
        (Section~\ref{sec:runningmd-protocols}).
  \item \textbf{Restraints:} \key{s} to keep the protocol's restraints
        (Section~\ref{sec:runningmd-restraints}).
  \item \textbf{Queue:} \key{k} to keep the order.
  \item \textbf{Hardware:} \key{3}, then \key{1} (sander) for every step
        (Section~\ref{sec:runningmd-engine}). Then \key{n}.
  \item \textbf{Review:} \key{3} (Write files and run locally now), then
        \key{1} and \key{y} (Section~\ref{sec:runningmd-local}).
\end{steps}

Several prompts in the setup guide, such as \prompt{Enter choice:}, have no
default. Pressing \key{Enter} alone gives \prompt{Received '', which is not
one of the available options}; type the letter or number.

\subsection{Step 1: Structure files}

ProPrep lists the topology--coordinate pairs it found in the project:

\begin{screen}
  Structure Pairs: 1 selected
    1. 1UBQ.prmtop + 1UBQ.rst7

  Auto-detected: 1 pairs available

What would you like to do?
   1. Use 1UBQ (exact match)
   2. Create custom pair (manual selection)
   3. Browse for structure files (topology + coordinates)
   4. Remove structure pair
\end{screen}

Use \key{2} to pick a topology and coordinates yourself, \key{3} to look in
another folder, and \key{4} to remove a pair you do not want. Several pairs
can be selected; each becomes its own simulation series. \key{h} explains the
pairing rules.

\begin{tip}
After you make a constant-pH file (\key{5a} \key{6},
Section~\ref{sec:topology-cpin}), ProPrep selects the \file{_cpin.prmtop}
pair instead (\prompt{Use 1UBQ (cpin-ready)}). For an ordinary simulation,
type \key{2} to add the normal pair and \key{4} to remove the cpin pair. In
the custom list the same file may appear twice; either entry works.
\end{tip}

\section{Protocols}\label{sec:runningmd-protocols}

\begin{screen}
Protocol Source:

   1. Use a predefined protocol     — standard multi-step protocols
   2. Load a saved protocol         — reuse a protocol from a previous project
   3. Build a custom protocol       — assemble steps from the template catalog
   4. Create a step from scratch    — launch the parameter wizard
\end{screen}

\subsection{Predefined protocols}

\begin{tabularx}{\linewidth}{@{}lL@{}}
\toprule
\textbf{Protocol} & \textbf{Stages} \\
\midrule
1 Membrane Protein Minimization (2) &
  Restrained, then unrestrained minimization, for membrane systems
  (Chapter~\ref{ch:membrane}). \\
2 Basic Equilibration (6) &
  Minimization; heating 0--300\,K (2\,ns); NPT density (1\,ns, CPU); NPT
  relaxation (3\,ns, GPU); NVT equilibration (4\,ns); NVT production
  (100\,ns). Based on Amber Tutorial~33. \\
3 Staged Equilibration 310\,K (9) &
  For metalloproteins and cofactors: minimization; heating to 310\,K; NPT
  density and relaxation; restraints released in stages (10, 5,
  2\,kcal/mol·Å²); 1\,ns unrestrained; 50\,ns NPT production. \\
4 Workshop Protocol (4) &
  1,000-cycle minimization; 20\,ps heating; 100\,ps NPT; 300\,ps NVT
  production. Made for laptops: about 75~minutes with \texttt{sander}. \\
5 York Lab Equilibration Tutorial (14) &
  A detailed 14-stage solvent-then-solute equilibration ending with 100\,ns
  production. \\
\bottomrule
\end{tabularx}

\medskip
After choosing a protocol you see its stages and these commands:

\begin{tabularx}{\linewidth}{@{}lL@{}}
\toprule
\textbf{Type} & \textbf{To} \\
\midrule
\typed{p 1} & Preview the mdin file of stage 1, with a comment on every line \\
\key{u} & Use the protocol as it is \\
\key{s} & Use only some stages (for example only the minimization) \\
\key{e} & Edit the protocol: change, remove, add or reorder stages \\
\key{b} & Back to the list \\
\bottomrule
\end{tabularx}

\subsection{Shortening or changing a protocol}\label{sec:runningmd-edit}

Type \key{e}, accept or type a protocol name, and the workflow editor opens.
Its commands:

\begin{tabularx}{\linewidth}{@{}lL@{}}
\toprule
\textbf{Command} & \textbf{What it does} \\
\midrule
\typed{edit 2} & Show stage 2's settings; then \typed{params} opens its mdin
  file in the text editor, \typed{name} and \typed{desc} rename it,
  \typed{files} changes its input and output names, \typed{back} returns \\
\typed{remove 3} & Remove stage 3 (asks for confirmation) \\
\typed{add before 2}, \typed{add after 3} & Insert a new stage \\
\typed{move 3 to 1} & Reorder stages \\
\typed{preview} & Show how files flow from stage to stage \\
\typed{info} & Change the protocol's name and description \\
\typed{save} & Save and continue the setup \\
\typed{cancel} & Discard the changes \\
\bottomrule
\end{tabularx}

\medskip
To make a stage shorter, change these lines in its mdin file:

\begin{itemize}
  \item \textbf{Minimization:} lower \texttt{maxcyc} (and \texttt{ncyc}).
        For ubiquitin, \texttt{ncyc=100, maxcyc=300} was enough.
  \item \textbf{Heating:} lower \texttt{nstlim}, and lower the step numbers
        in the \texttt{\&wt} lines (\texttt{ISTEP1}, \texttt{ISTEP2}) to match,
        or the temperature ramp will not fit the new length.
  \item \textbf{Production:} lower \texttt{nstlim}, and lower \texttt{ntwx}
        and \texttt{ntpr} so you still get enough snapshots. With
        \texttt{nstlim=500, ntwx=25} you get 20 snapshots.
\end{itemize}

In the text editor (\texttt{nano}), move with the arrow keys, type your
change, then press \key{Ctrl-X}, \key{y} and \key{Enter}. ProPrep prints
\prompt{MDIN parameters updated} and lists the parameters again.

\begin{watchout}
Some screens of the editor do not match what is really run in version
1.23.0:
\begin{itemize}
  \item ProPrep does not check your edits. A minimization saved without its
        \texttt{imin=1} line was accepted without a warning. Read the
        parameter list after every edit.
  \item \typed{preview} shows \texttt{pmemd} commands and file names such as
        \file{step_1.rst}, whatever engine you choose later. The real run uses
        your engine and its own file names. Changing the names with
        \typed{files} can break the preview's chain; removing a stage renumbers
        the chain correctly on its own, so leave \typed{files} alone.
  \item \typed{info} does not show information: it asks for a new name and
        description.
  \item \typed{preview} ends with \prompt{Preview template content for a
        step? [y/n]}, which accepts only \key{y} or \key{n}.
  \item The Review table at the end of setup shows the protocol's
        \emph{original} lengths, not your shortened ones. The files written
        do contain your changes.
\end{itemize}
\end{watchout}

\subsection{Building a protocol from the catalog}

Choice \key{3} (Build a custom protocol) lists all 35 built-in stages from
the five protocols, numbered 1--35. Type numbers to add stages (\typed{18,21}),
\typed{p 18} to preview one, \typed{c 18 22} to compare two side by side,
\key{r} to remove the last one, and \key{d} when done; then give the protocol
a name. Templates you imported yourself are not listed in this catalog.

\subsection{A stage from scratch}

Choice \key{4} starts a wizard that asks about every Amber setting in 13
sections, with explanations at three levels of detail (Terse, Compact,
Verbose). It covers minimization, dynamics and analysis runs. We made a
minimization this way and ran it successfully with \texttt{sander}.

\begin{watchout}
With all defaults accepted, the wizard's minimization contained
constant-pressure settings (\texttt{ntb=2}, \texttt{ntp=1},
\texttt{barostat=2}), which do not belong in a minimization. \texttt{sander}
ran it, with the note \prompt{mcbarint is greater than the number of steps.
This is effectively constant volume.} Answer the pressure questions with care,
or use a predefined minimization stage instead.
\end{watchout}

\subsection{Saved protocols and the template library}

When you \typed{save} an edited protocol, ProPrep stores it in
\file{~/.proprep/md_workflows/user/custom/} in your home folder, which is
shared by all your projects. The library (\key{5b} \key{5}) lets you import
mdin templates and protocol files, see statistics (35 built-in templates:
8 minimization, 5 heating, 10 equilibration, 4 production, 8 other), preview
templates, and delete your own templates and protocols.

\begin{watchout}
In version 1.23.0:
\begin{itemize}
  \item \key{2} (Load a saved protocol) always said \prompt{No saved protocols
        found}, even when the library listed one. To reuse a protocol, use the
        import option there with an exported \file{.json} file.
  \item Deleting custom protocols in the library reports success, but the
        files stay in \file{~/.proprep/md_workflows/user/custom/}. Delete them
        by hand if needed.
  \item An imported \file{.mdin} template is saved in the project
        (\file{md_templates/user/custom/}), not in the folder the statistics
        screen names.
  \item The statistics show \prompt{? steps} for every custom protocol.
\end{itemize}
\end{watchout}

\section{Restraints}\label{sec:runningmd-restraints}

Step~3 of the setup shows the restraints of each stage:

\begin{screen}
Protocol Restraints:

    #    Step                   Mask                 Weight    DISANG
    1    Energy Minimization    @CA,C,O,N&!:WAT        10.0    —
    2    System Heating         @CA,C,O,N&!:WAT         1.0    —
    3    NVT Production MD      (none)                    —    —

Options:
   (e) Edit positional restraints
   (g) Generate new restraint mask (atom type, chain, redox site selection)
   (d) Configure DISANG restraints (distance/angle/torsion)
   (r) Configure GROUP restraints (advanced: multi-group with FIND criteria)
   (s) Continue with current restraints
\end{screen}

The predefined protocols already restrain sensible atoms. Type \key{s} unless
you have a reason to change them.

\subsection{Positional restraints and masks}

A \term{positional restraint} pulls atoms back towards their starting
positions, like a spring. The \emph{weight} is the spring strength in
kcal/mol·Å²: 10 is firm, 1 is gentle, 0.1 is very light.

The atoms are chosen with an Amber \term{mask}, a short code:

\begin{tabularx}{\linewidth}{@{}lL@{}}
\toprule
\textbf{Mask} & \textbf{Means} \\
\midrule
\texttt{:1-76} & Residues 1 to 76 (\texttt{:} = residue numbers) \\
\texttt{@CA} & All atoms named CA (\texttt{@} = atom names) \\
\texttt{@CA,C,O,N} & The backbone atoms \\
\texttt{!@H=} & Everything except hydrogens (\texttt{!} = not, \texttt{=} = any ending) \\
\texttt{\&} & ``and'': \texttt{@CA\&:1-70} is the CA atoms of residues 1--70 \\
\texttt{!:WAT} & Not water \\
\bottomrule
\end{tabularx}

\medskip
\key{e} lets you choose stages (\typed{2}, \typed{2-4} or \typed{all}), type a
mask (or \typed{none} to remove restraints) and set a weight for each stage.
\key{g} opens a mask generator that builds the mask from answers:

\begin{screen}
═══ General Structural Restraints ═══
Choose atom types to restrain:
  (b) Backbone atoms (@CA,C,O,N)
  (h) Heavy atoms only (!@H=)
  (c) Alpha carbons only (@CA)
  (u) Custom atom selection (you specify)
Select atom type (c): §ans«c»
...
Chain selection (all): §ans«A:1-70»
...
Exclude water molecules from restraints? [y/n] (y): §ans«y»
...
│ FINAL COMBINED MASK (saved to workspace):                │
│   @CA&:1-70&!:WAT                                        │
\end{screen}

After \key{a} (accept), ProPrep asks which stages should use the new mask and
with what weight. In our test the heating stage received
\texttt{@CA\&:1-70\&!:WAT} with weight 2.0, and the run used it.

\begin{tip}
Residue numbers in masks count from 1 through the whole system, including
ions. For ubiquitin the generator reported \prompt{Chain A: residues 1-102}:
76 amino acids plus 26 ions.
\end{tip}

\subsection{Distance restraints (DISANG)}

A \term{DISANG} file holds restraints on distances, angles or torsions, for
example ``keep these two atoms between 35 and 40\,Å apart''. \key{d} opens
the restraint manager:

\begin{screen}
═══ MD RESTRAINT CONFIGURATION ═══
DISANG Restraints (Distance/Angle/Torsion)
  1. ✓ Import DISANG restraints from file
  2. ✓ Configure DISANG restraints interactively
  3. ○ Display current DISANG restraints
  4. ○ Export restraints to DISANG file
...
Choose option (or 'd' when done, 'x' to exit) (2):
\end{screen}

The route that worked:

\begin{steps}
  \item Write the restraint file yourself and give it a name ending in
        \file{.disang} in the project folder. ProPrep lists only
        \file{*.disang} files there.
  \item Type \key{1} and accept the suggested file.
  \item Type \key{3} to check that the restraint is listed.
  \item Type \key{4} to export it (default \file{restraints.disang}).
  \item Type \key{d}, choose the stages (for example \typed{3}) and the
        \prompt{Restraint dump frequency} (how often the distances are
        written).
\end{steps}

\begin{watchout}
In version 1.23.0:
\begin{itemize}
  \item Importing printed \prompt{Error importing restraints:
        'MDRestraintModule' object has no attribute 'get\_from\_workspace'},
        yet the restraint was imported and later used correctly by
        \texttt{sander}. Check with \key{3} and carry on.
  \item Interactive entry (\key{2}) could not find any residue (\prompt{Residue
        A:1 not found}): the structure ProPrep extracts from the topology has
        no chain letters. Use a file instead.
\end{itemize}
\end{watchout}

For a stage with DISANG restraints, ProPrep adds \texttt{nmropt=1} and
three lines to its mdin file:
\texttt{DISANG=}, \texttt{DUMPAVE=} and \texttt{LISTOUT=}. The measured
distances appear in \file{restraints_dump.txt} in the stage folder.

\subsection{GROUP restraints}

\key{r} sets up \term{GROUP restraints}, an older Amber format that restrains
several groups of atoms with different strengths. ProPrep asks for the number
of groups (1--10), a title and force constant for each, search criteria
(\texttt{CA * * *} means every atom named CA) and up to seven residue ranges,
then shows the block and asks which stages should use it. They cannot be
combined with a positional mask in the same stage. We configured a group but
did not apply it to a run.

\section{Choosing the engine and hardware}\label{sec:runningmd-engine}

The \term{MD engine} is the program that does the calculation. Amber has
several:

\begin{tabularx}{\linewidth}{@{}lL@{}}
\toprule
\textbf{Engine} & \textbf{What it is} \\
\midrule
\texttt{sander} & AmberTools engine, one processor core. Installed with
  ProPrep. \\
\texttt{sander.MPI} & The same, shared over several cores with \term{MPI}
  (a system for splitting work between processors). Installed with ProPrep. \\
\texttt{pmemd} & Faster engine from the separate Amber package. \textbf{Not
  installed with ProPrep.} \\
\texttt{pmemd.MPI} & \texttt{pmemd} over several cores. Not installed. \\
\texttt{pmemd.cuda} & \texttt{pmemd} on an NVIDIA graphics card (GPU), by far
  the fastest. Not installed. \\
\bottomrule
\end{tabularx}

\begin{needsoftware}
\texttt{pmemd}, \texttt{pmemd.MPI} and \texttt{pmemd.cuda} are part of Amber,
which is obtained and installed separately from AmberTools, and
\texttt{pmemd.cuda} also needs an NVIDIA GPU. Without them, use
\texttt{sander} or \texttt{sander.MPI}: everything in this chapter works, only
more slowly. ProPrep can still write input files and cluster scripts for
\texttt{pmemd} to run on another computer.
\end{needsoftware}

Step~5 of the setup:

\begin{screen}
  Minimization on CPU: GPU fixed-point precision (SPFP) can silently clip
  large forces from atomic clashes in unrelaxed starting structures.
  NPT equilibration on CPU: the GPU PME grid is fixed at run start and
  cannot reorganize as box dimensions change during density equilibration.
  NVT/production on GPU: once the structure is relaxed and density has
  converged, GPU acceleration is safe and dramatically faster.

  Available: 11 CPUs, 0 GPUs

  Simulations: 3 configured
    3x pmemd

What would you like to do?
   1. Apply recommended assignments (CPU for min/NPT, GPU for rest)
   2. Configure individual simulation engines
   3. Set all simulations to same engine
   p. Load cluster profile
   l. Load run plan (requires cluster profile)
   s. Save current assignments as run plan (requires cluster profile)
\end{screen}

\begin{itemize}
  \item \key{3} sets one engine for every stage. This is the simplest way to
        choose \texttt{sander}.
  \item \key{2} goes through the stages one by one. For \texttt{sander.MPI}
        it also asks \prompt{Number of MPI processes}; for
        \texttt{pmemd.cuda}, which GPU to use. This list marks missing
        engines \prompt{(not installed)}.
  \item \key{1} applies ProPrep's recommendation: CPU engines for
        minimization and NPT, GPU for the rest.
  \item \key{h} explains why some stages should not run on a GPU.
\end{itemize}

\begin{watchout}
In version 1.23.0 every stage starts with the engine \texttt{pmemd}, which is
not installed with ProPrep. The recommended assignment (\key{1}, also the
default) chooses \texttt{pmemd.cuda} for heating and production even when the
screen says \prompt{0 GPUs}, and only prints \prompt{No GPUs detected on this
machine. Defaulting to GPU 0.} The list under \key{3} does not mark missing
engines. If you run with a missing engine, the run stops with
\prompt{Error starting simulation: [Errno 2] No such file or directory:
'pmemd'} and \prompt{Simulations started: 0}. Go back to \key{5b} \key{1}
and choose \texttt{sander} at this step; the queue menu of the Execute screen
cannot change the engine.
\end{watchout}

How many processes? For ubiquitin (15,729 atoms) on an 11-core laptop,
\texttt{sander.MPI} with 4 processes ran 2\,ps of heating in 22\,s, about
7.7\,ns per day. Leave some cores free for the rest of your computer.

Cluster profiles (\key{p}) and run plans (\key{l}, \key{s}) are for
supercomputers; see Section~\ref{sec:runningmd-cluster}.

\section{Running on your computer}\label{sec:runningmd-local}

Step~6 shows a summary and the ways to finish:

\begin{screen}
What would you like to do?
   1. Save setup to workspace (execute later with 'run' command)
   2. Write MDIN only (no SLURM, no launch)
   3. Write files and run locally now
   4. Write MDIN + SLURM scripts (no launch) (manual wizard — no profile loaded)
   5. Write files and submit via sbatch now (requires cluster profile)
\end{screen}

\begin{itemize}
  \item \key{3} writes the files and opens the Execute Simulations menu. Type
        \key{1} (Execute simulation queue) and \key{y}.
  \item \key{2} only writes the input files and a script
        \file{run_workflow.sh} that runs all stages in a row. You can start
        the run later from \key{5b} \key{2}, or copy the folder to another
        computer. ProPrep also asks \prompt{Add extended production run
        cycles?}: \key{y} and a number adds that many extra production runs,
        each continuing from the last.
  \item \key{4} and \key{5} are for clusters
        (Section~\ref{sec:runningmd-cluster}).
  \item \key{1} saves the setup without writing files (not tried for this
        manual).
\end{itemize}

When the run starts, ProPrep launches the first stage in the background and
starts each following stage as soon as the one before has finished:

\begin{screen}
Executing simulation 1/3: 1UBQ_1UBQ_Energy Minimization
...
✓ Simulation '1UBQ_1UBQ_Energy Minimization' started in background (PID: 4080)
...
Executing simulation 2/3: 1UBQ_1UBQ_System Heating
  Queued (depends on: 1UBQ_1UBQ_Energy Minimization)
Will auto-start when dependency completes
...
===== Execute Simulations =====

▶ Running: Energy Minimization
   Next: System Heating (+1 more pending)
\end{screen}

You can keep using ProPrep while it runs. The Execute menu shows
\prompt{Finished} for each completed stage.

\subsection{Where the files go}

Each run makes a new folder \file{simulations/batch_<date>_<time>/} in the
project. (Choosing \key{2} and then running from \key{5b} \key{2} gives two
batch folders: one with the written files, one with the actual run.)

\begin{tabularx}{\linewidth}{@{}lL@{}}
\toprule
\textbf{In each stage folder} & \textbf{Contents} \\
\midrule
\file{simulation.mdin} & The settings for this stage \\
\file{run_simulation.sh} & The exact command that was run \\
\file{*.mdout} & Energies, temperature and pressure through the stage \\
\file{*.rst7} & Final coordinates (the start of the next stage) \\
\file{*.nc} & The trajectory (dynamics stages only) \\
\file{mdinfo} & Progress of the running stage \\
\file{simulation.stdout}, \file{.stderr} & Messages from the engine; look here if a stage fails \\
\file{config.json} & ProPrep's record of the stage \\
\bottomrule
\end{tabularx}

\medskip
The \file{shared/} folder holds the topology and starting coordinates. Inside
each stage folder the incoming coordinates are always called
\file{<name>.rst7}, even when they are really the previous stage's result.
The chaining is correct; only the name is the same.

\subsection{How long it takes}

Measured for ubiquitin (15,729 atoms) on an 11-core Mac laptop:

\begin{tabularx}{\linewidth}{@{}lllL@{}}
\toprule
\textbf{Stage} & \textbf{Engine} & \textbf{Length} & \textbf{Time} \\
\midrule
Minimization & sander.MPI, 4 processes & 300 cycles & 6\,s \\
Heating & sander.MPI, 4 processes & 1,000 steps (2\,ps) & 22\,s \\
Production & sander & 500 steps (1\,ps) & 33\,s \\
\bottomrule
\end{tabularx}

\medskip
At this rate the full Workshop protocol (420\,ps) takes over an hour, and
100\,ns would take weeks. Long simulations belong on a computer with a GPU and
\texttt{pmemd.cuda}, usually a cluster.

\section{Watching a running simulation}\label{sec:runningmd-monitor}

Choose \key{3} in the Execute menu (or \key{5b} \key{3}) and pick a
simulation:

\begin{screen}
AMBER Simulation Monitor
Data & Status:
 1. Show current status
 2. Display energy summary
 3. Display temperature summary
 4. Display pressure summary
 5. Show timing information

ASCII Plots:
 6. Plot energy (total)
 7. Plot temperature
 8. Plot pressure
 9. Plot kinetic energy
10. Plot potential energy
11. Plot RMS gradient

Periodic box:
12. Distance of the solute to its periodic images, so far

Control:
 r. Refresh data
 b. ← Back
\end{screen}

\key{1} shows the current step, time, temperature and energy. \key{6}--\key{11}
draw simple text plots in the window. \key{r} reads the latest values.

Option \key{12} measures, for every snapshot so far, the shortest distance
between the protein and its own copy in the neighboring box. It should stay
well above the cutoff distance (usually 10\,Å); otherwise the protein can
feel itself across the box edge and the box is too small. For ubiquitin's
heating stage:

\begin{screen}
│ Box                       │ truncated octahedron; edges a 61.99-61.99, ...
│ Closest approach          │ 29.28 Å at frame 17                         │
│ Mean                      │ 30.01 ± 0.35 Å                              │
│ Nonbonded cutoff          │ 10 Å                                        │
│ Frames below the cutoff   │ 0 of 20                                     │

The solute never comes within the 10 Å cutoff of its own image in the frames
analysed; at its closest it is 19.28 Å beyond it.
\end{screen}

\prompt{Show timing information} said \prompt{No timing information
available} for a stage that had just finished.

\section{Running on a computing cluster}\label{sec:runningmd-cluster}

Most university clusters use a job system called \term{SLURM}: you submit a
script, and it runs when processors are free. ProPrep can write these scripts.
We wrote them for this manual but had no cluster to run them on.

At the Review step, \key{4} (Write MDIN + SLURM scripts) writes the input
files and then asks for the job settings:

\begin{tabularx}{\linewidth}{@{}lL@{}}
\toprule
\textbf{Question} & \textbf{Default} \\
\midrule
Job name; partition (queue) & \texttt{amber\_simulation}; \texttt{gpu} \\
Nodes; tasks per node; CPUs per task & 1; 1; 4 \\
Use GPU?; number; type & yes; 1; any \\
Memory; time limit & 16G; 02:00:00 \\
Email notifications & no \\
Load modules? & yes, \texttt{amber/22} \\
Conda environment; advanced options & no; no \\
\bottomrule
\end{tabularx}

\medskip
Ask your cluster's support team for the right partition, modules and
account. The scripts go to the folder \file{simulations/slurm_scripts/}.
There is one \file{stepN.sh} per stage, plus \file{submit_workflow.sh},
which submits them all so that each waits for the one before. On the cluster you would run
\texttt{./submit\_workflow.sh} and watch with \texttt{squeue -u \$USER}.

\begin{watchout}
Check the scripts before you submit them. In our test (version 1.23.0):
\begin{itemize}
  \item With all defaults, a GPU was requested on a \texttt{gpu} partition
        even though the engine was \texttt{sander}, which cannot use a GPU.
  \item For a constant-pH production stage, \file{step2.sh} used the
        ordinary topology instead of the \file{_cpin.prmtop} one. The local
        \file{run_workflow.sh} was correct.
  \item \file{step2.sh} contains a line \texttt{\#SBATCH
        --dependency=afterok:\$STEP1\_JOBID}, which fails if you submit that
        script on its own. Use \file{submit_workflow.sh}.
\end{itemize}
\end{watchout}

\subsection{Cluster profiles and run plans}

A \term{cluster profile} stores the details of one cluster (job system,
modules, partitions, which engine to use on CPU or GPU nodes) so you do not
type them each time. ProPrep includes one example, \textsf{TAMU ACES}. Manage
profiles with \key{5b} \key{6}: show, fill in required fields (such as your
account), edit the full file, create a new profile from a template, export,
import and delete. A \term{run plan} remembers which engine and resources
each stage of a protocol uses; it is saved at the hardware step (\key{s}) and
managed with \key{5b} \key{7}. Review option~\key{5} (submit with
\texttt{sbatch} now) needs a loaded profile and a SLURM system; it was not
tried for this manual.

\section{Constant-pH simulations}\label{sec:runningmd-cph}

If the project has a cpin file (Section~\ref{sec:topology-cpin}), ProPrep
offers constant-pH MD after you choose the protocol:

\begin{screen}
Constant pH MD Available
  CPIN file: 1UBQ.cpin (8 residues)
  Solvent model: explicit
  Modified topology: 1UBQ_cpin.prmtop
...
Enable Constant pH MD for production steps? [y/n] (y): §ans«y»
Target pH (7.0): §ans«»
Steps between protonation state attempts (ntcnstph) (100): §ans«»
Solvent relaxation steps after protonation change (ntrelax) (200): §ans«»
Salt concentration in M (saltcon) (0.1): §ans«»
✓ Constant pH MD enabled
  icnstph=2, solvph=7.0, ntcnstph=100, ntrelax=200, saltcon=0.1
  -cpin -cpout -cprestrt will be applied to production steps only
\end{screen}

The \prompt{saltcon} value is not the salt in your water box; leave it at
0.1. Constant-pH settings apply only to production stages. The default is
\key{y}; answer \key{n} for an ordinary simulation. We wrote constant-pH
input files and scripts but did not run a constant-pH simulation for this
manual.

\section{If something goes wrong}\label{sec:runningmd-problems}

\begin{tabularx}{\linewidth}{@{}lL@{}}
\toprule
\textbf{What you see} & \textbf{Why, and what to do} \\
\midrule
\prompt{No such file or directory: 'pmemd'}, \prompt{Simulations started: 0} &
  The engine is \texttt{pmemd} (the default), which is not installed. Redo
  \key{5b} \key{1} and choose \texttt{sander} at the hardware step. \\
\prompt{Received '', which is not one of the available options} &
  That prompt has no default. Type a letter or number. \\
The Review table shows the old lengths &
  Display only; your edited values are in the files. Open
  \file{simulation.mdin} in a stage folder to confirm. \\
\prompt{No saved protocols found} &
  Known limitation; import the protocol's \file{.json} file instead. \\
\prompt{Residue A:1 not found} in DISANG entry &
  Use a \file{.disang} file and import it. \\
A stage stops early or never finishes &
  Read \file{simulation.stderr} and the end of the \file{.mdout} file in the
  stage folder. \\
Wrong Amber programs are used &
  Another Amber installation is active (\texttt{AMBERHOME}); see
  Section~\ref{sec:tutorial-before}. \\
\bottomrule
\end{tabularx}

\medskip
See also Chapters~\ref{ch:troubleshooting} and~\ref{ch:knownissues}.
